%%%%%%%%%%%%%%%%%%%%%%%%%%%%%%%%%%%%%%%%%%%%%%%%%%%%%%%%%%%%%%%%%%%%%%
% Lớp 11 | Giới hạn. Hàm số liên tục | Giới hạn của hàm số
% Dạng 1D3?2-1: Câu hỏi lý thuyết
% Số câu: 35  (sinh tự động bởi tach_cau.py; từ nay thầy sửa trực tiếp file này)
%%%%%%%%%%%%%%%%%%%%%%%%%%%%%%%%%%%%%%%%%%%%%%%%%%%%%%%%%%%%%%%%%%%%%%

%%%%%%%%%%%%%%%% TRẮC NGHIỆM 4 PHƯƠNG ÁN %%%%%%%%%%%%%%%%

% nguồn: Ok_TracNghiemLop11_Lan1
\begin{ex}%[1D3N2-1]%[ma:411912e7]
	Giá trị của $\lim\limits_{x\to+\infty}x^3$  bằng
	\choice 
		 { $-\infty $}
		{\True $+\infty $} 
		 { $0$}
		{  $3$}
	 \loigiai{
		Do $3$ là số nguyên dương nên  $\lim\limits_{x\to+\infty}x^3=+\infty$.
	} 
\end{ex}

% nguồn: Ok_TracNghiemLop11_Lan1
\begin{ex}%[1D3N2-1]%[ma:c84f45ce]
	Cho $\lim\limits_{x\to 2}f(x)=3$. Giá trị của $\lim\limits_{x \to  2}\left[f(x)+x\right]$ bằng
	 \choice
		 {\True $5$}
		 {$6$}
		 {$1$}
		 {$4$}
	\loigiai{ 
	Ta có $\lim\limits_{x \to 2}\left[f(x)+x\right] =  \lim\limits_{x \to 2}f(x) + \lim\limits_{x \to 2}x = 3+2=5$.
	 }
\end{ex}

% nguồn: Ok_TracNghiemLop11_Lan1
\begin{ex}%[1D3N2-1]%[ma:c6fbcd70]
	Chọn khẳng định  \textbf{sai}.
	 \choice
		{Hàm  số đa thức liên tục trên $\mathbb{R}$}
		{\True Hàm  số $y=f(x)$ liên tục trên đoạn $[a; b]$ nếu nó liên tục trên khoảng $(a; b)$}
		{Hàm số $y=f(x)$ liên  tục tại điểm $x_0$ nếu $\lim\limits_{x\to x_0} f(x)=f\left(x_0\right)$}
		{Hàm số $y=f(x)$ liên tục trên một khoảng nếu  nó liên tục tại mọi điểm của khoảng đó}
	\loigiai{Hàm số $y=f(x)$ liên tục trên  đoạn $[a; b]$ nếu nó liên tục trên khoảng $(a; b)$ và \\
		$\lim\limits_{x\to a^{+}} f(x)=f(a),  \lim\limits_{x\to a^{-}} f(x)=f(b)$.
	 }
\end{ex}

% nguồn: Ok_TracNghiemLop11_Lan1
\begin{ex}%[1D3N2-1]%[ma:a10c99e5]
	Cho hàm số $y=f(x)$ xác định  trên $(a;b)$ và $x_0\in (a;b)$. Giả sử các giới hạn (hữu hạn) sau tồn tại, giới hạn nào là đạo hàm của hàm số $y=f(x)$ tại điểm $x_0$ ?
	\choice 
		{\True  $\lim\limits_{x\to x_0} \dfrac{f(x)-f(x_0)}{x-x_0}$}
		{$\lim\limits_{\Delta x\to +\infty} \dfrac{\Delta y}{\Delta  x}$}
		{$\lim\limits_{\Delta x\to x_0} \dfrac{\Delta y}{\Delta  x}$}
		{$\lim\limits_{x\to 0} \dfrac{f(x)-f(x_0)}{x-x_0}$} 
	\loigiai{ 
	Theo định nghĩa đạo hàm của hàm số  tại một điểm.}
\end{ex}

% nguồn: Ok_Toan11_TracNghiem_Lan2
\begin{ex}%[1D3H2-1]%[ma:e41067ad]
	Cho hàm số $y=f(x)$ có giới hạn hữu  hạn tại $x=1$ và thỏa mãn $\lim\limits_{x\to 1}\dfrac{f(x)-1}{x-1}$ tồn tại hữu hạn. Tính $\lim\limits_{x\to 1} f(x)$.
	\choice 
		{$\lim\limits_{x\to 1} f(x)=0$} 
		{\True  $\lim\limits_{x\to 1} f(x)=1$}
		{$\lim\limits_{x\to 1}  f(x)=7$}
		{$\lim\limits_{x\to  1} f(x)=+\infty$}
	 \loigiai{
		Nếu $\lim\limits_{x\to 1} f(x)\neq 1$ thì $\lim\limits_{x\to  1}\left[f(x)-1\right]\neq 0$.\\
		Hơn nữa, $\lim\limits_{x\to 1} (x-1)=0$ nên $\lim\limits_{x\to 1}\dfrac{f(x)-1}{x-1}$ không thể  hữu hạn.\\
		 Vậy $\lim\limits_{x\to 1} f(x)=1$.
	} 
\end{ex}

% nguồn: Ok_Toan11_TracNghiem_Lan2
\begin{ex}%[1D3H2-1]%[ma:8cd6b4d8]
	Cho hàm số $y=f(x)$ có giới hạn  hữu hạn tại $x=2$ và thỏa mãn $\lim\limits_{x\to 2}\dfrac{3f(x)-7}{2x-4}$ tồn tại hữu hạn. Tính $\lim\limits_{x\to 2} f(x)$.
	\choice 
		{\True  $\lim\limits_{x\to 2} f(x)=\dfrac{7}{3}$}
		 {$\lim\limits_{x\to 2} f(x)=\dfrac{14}{3}$}
		{$\lim\limits_{x\to 2}  f(x)=7$}
		{$\lim\limits_{x\to  2} f(x)=5$}
	 \loigiai{
		Nếu $\lim\limits_{x\to 2} f(x)\neq\dfrac{7}{3}$ thì $\lim\limits_{x\to 2}\left[3f(x)-7\right]\neq 0$.\\ 
		Hơn nữa, $\lim\limits_{x\to 2} (2x-4)=0$ nên $\lim\limits_{x\to 2}\dfrac{3f(x)-7}{2x-4}$ không thể hữu hạn.\\ 
		Vậy $\lim\limits_{x\to  2} f(x)=\dfrac{7}{3}$.
	} 
\end{ex}

% nguồn: Ok_Toan11_TracNghiem_Lan2
\begin{ex}%[1D3V2-1]%[ma:c379f4de]
	Cho hàm số $y=f(x)$ có giới hạn hữu hạn tại $x=-1$ và thỏa mãn $\lim\limits_{x\to-1}\dfrac{2 f(x)-1}{x+1}=3$. Hỏi giá trị của $\lim\limits_{x\to-1}\dfrac{\sqrt{4 f(x)+7}-3}{x^2-1}$ thuộc khoảng nào  sau đây?
	\choice 
		{\True  $(-1;0)$}
		 {$(-2;-1)$}
		 {$(0;1)$}
		{$(1;2)$} 
	\loigiai{ 
		Nếu $\lim\limits_{x\to-1}[2 f(x)] \neq 1$ thì $\lim\limits_{x\to-1}[2 f(x)-1] \neq 0$ và $\lim\limits_{x\to-1}(x+1)=0$.\\ 
		 Khi đó $\lim\limits_{x\to-1}\dfrac{2 f(x)-1}{x+1}$ không thể hữu hạn.\\
		Do đó,  $\lim\limits_{x\to-1}[2 f(x)]=1 \Rightarrow \lim\limits_{x\to-1}f(x)=\dfrac{1}{2}$.\\
		Ta có 
		\allowdisplaybreaks 
		\begin{eqnarray*}
			 L&=& \lim\limits_{x\to-1}\dfrac{\sqrt{4 f(x)+7}-3}{x^2-1}\\
			&=&  \lim\limits_{x\to-1}\dfrac{4 f(x)-2}{(x+1)(x-1)\left(\sqrt{4f(x)+7}+3\right)}\\
			&=&  \lim\limits_{x\to-1}\left[\dfrac{2f(x)-1}{x+1}\cdot \dfrac{2}{(x-1)\left(\sqrt{4f(x)+7}+3\right)}\right]\\
			&=&  3\cdot \dfrac{2}{-2(3+3)}\\
			&=& -\dfrac{1}{2}\in  (-1;0).
		\end{eqnarray*}
	 }
\end{ex}

% nguồn: Ok_Toan11_TracNghiem_Lan2
\begin{ex}%[1D3V2-1]%[ma:820bfaf8]
	Cho hàm số $y=f(x)$ có giới hạn hữu hạn tại $x=0$ và  thỏa mãn $\lim\limits_{x\to 0}\dfrac{2 f(x)+1}{x^2}$ tồn tại hữu hạn. Tính $\lim\limits_{x\to 0}f(x)$.
	\choice 
		 {$\lim\limits_{x\to 0}f(x)=-1$}
		 {$\lim\limits_{x\to 0}f(x)=1$}
		 {$\lim\limits_{x\to 0}f(x)=\dfrac{1}{2}$}
		{\True  $\lim\limits_{x\to 0}f(x)=-\dfrac{1}{2}$}
	\loigiai{ 
		Nếu $\lim\limits_{x\to 0}[2f(x)] \neq-1$ thì $\lim\limits_{x\to 0}[2 f(x)+1] \neq 0$ và  $\lim\limits_{x\to 0}x^2=0$.\\
		Khi đó $\lim\limits_{x\to 0}\dfrac{2 f(x)+1}{x^2}$ không  thể hữu hạn.\\
		 Do đó, $\lim\limits_{x\to 0}[2 f(x)]=-1 \Rightarrow \lim\limits_{x\to 0}f(x)=-\dfrac{1}{2}$
	 }
\end{ex}

% nguồn: Ok_Toan11_TracNghiem_Lan2
\begin{ex}%[1D3V2-1]%[ma:e099d7aa]
	Cho hàm số $y=f(x)$ có giới hạn hữu hạn tại $x=1$ và thỏa mãn $\lim\limits_{x\to 1}\dfrac{f(x)-5}{x-1}=8$.  Tính $L=\lim\limits_{x\to 1}\dfrac{\sqrt{f(x)+4}-3}{x-1}$.
	 \choice
		 {$L=8$}
		 {\True $L=\dfrac{4}{3}$}
		 {$L=40$}
		{$L=-40$} 
	\loigiai{ 
		Nếu $\lim\limits_{x\to 1}f(x) \neq 5$ thì $\lim\limits_{x\to 1}[f(x)-5] \neq 0$ và $\lim\limits_{x\to  1}(x-1)=0$.\\
		Khi  đó $\lim\limits_{x\to 1}\dfrac{f(x)-5}{x-1}$ không thể hữu hạn. Do đó, $\lim\limits_{x\to 1}f(x)=5$.\\
		Ta  có
		 \allowdisplaybreaks
		\begin{eqnarray*}
			L&=& \lim\limits_{x\to 1}\dfrac{\sqrt{f(x)+4}-3}{x-1}\\ 
			&=&  \lim\limits_{x\to 1}\dfrac{f(x)-5}{(x-1)\left(\sqrt{f(x)+4}+3\right)}\\
			&=& \lim\limits_{x\to 1}  \left[\dfrac{f(x)-5}{x-1}\cdot \dfrac{1}{\sqrt{f(x)+4}+3}\right]\\
			&=& 8\cdot \dfrac{1}{3+3}=\dfrac{4}{3}. 
		\end{eqnarray*}
	} 
\end{ex}

% nguồn: Ok_Toan11_TracNghiem_Lan2
\begin{ex}%[1D3V2-1]%[ma:4068fe0f]
	Cho  hàm số $y=f(x)$ có giới hạn hữu hạn tại $x=0$ và thỏa mãn $\lim\limits_{x\to 0}\dfrac{3 f(x)-12}{x}=7$. Biết $\lim\limits_{x\to 0}\dfrac{\sqrt{f(x)}-2}{x}=\dfrac{a}{b}$, ($a$, $b \in \mathbb{Z}$, $b>0$) với $\dfrac{a}{b}$ là phân số tối giản. Tính $T=a+b$.
	\choice 
		 {$T=18$}
		{\True  $T=19$}
		{$T=20$} 
		{$T=21$} 
	 \loigiai{
		 Nếu $\lim\limits_{x\to 0}[3f(x)] \neq 12$ thì $\lim\limits_{x\to 0}[3f(x)-12] \neq 0$ và $\lim\limits_{x\to 0}x=0$.\\
		Khi  đó $\lim\limits_{x\to 0}\dfrac{3 f(x)-12}{x}$ không thể hữu hạn.\\
		Do đó, $\lim\limits_{x\to 0}[3f(x)]=12  \Rightarrow \lim\limits_{x\to 0}f(x)=4$.\\
		 Ta có
		 \allowdisplaybreaks
		\begin{eqnarray*}
			L&=&\lim\limits_{x\to  0}\dfrac{\sqrt{f(x)}-2}{x}\\
			 &=&\lim\limits_{x\to 0}\dfrac{f(x)-4}{x\left(\sqrt{f(x)}+2\right)}\\
			 &=&\lim\limits_{x\to 0 }\left[\dfrac{3f(x)-12}{x}\cdot \dfrac{1}{3\left(\sqrt{f(x)}+2\right)}\right]\\
			 &=&7\cdot \dfrac{1}{3(2+2)}=\dfrac{7}{12}.
		\end{eqnarray*}
		Vậy $a=7$; $b=12$ và  $T=a+b=7+12=19$.
	 }
\end{ex}

% nguồn: Ok_Toan11_TracNghiem_Lan2
\begin{ex}%[1D3V2-1]%[ma:a1b0e0c2]
	Cho hàm số $y=f(x)$ có giới hạn hữu hạn tại $x=4$ và thỏa mãn $\lim\limits_{x\to 4}\dfrac{f(x)+3}{x-4}=1$. Tính  $L=\lim\limits_{x\to 4}\dfrac{\sqrt[3]{f(x)+4}-1}{x-4}$.
	\choice 
		 {$L=\dfrac{1}{4}$}
		{\True $L=\dfrac{1}{3}$} 
		 {$L=\dfrac{2}{3}$}
		{$L=\dfrac{1}{2}$} 
	\loigiai{ 
		Nếu  $\lim\limits_{x\to 4}f(x) \neq-3$ thì $\lim\limits_{x\to 4}[f(x)+3] \neq 0$ và $\lim\limits_{x\to 4}(x-4)=0$.\\
		Khi  đó $\lim\limits_{x\to 4}\dfrac{f(x)+3}{x-4}$ không thể hữu hạn.\\
		Do đó, $\lim\limits_{x\to  4}f(x)=-3$.\\
		 Ta có
		 \allowdisplaybreaks
		\begin{eqnarray*}
			L&=&\lim\limits_{x\to  4}\dfrac{\sqrt[3]{f(x)+4}-1}{x-4}\\
			 &=&\lim\limits_{x\to 4}\dfrac{f(x)+3}{(x-4)\left[\sqrt[3]{\left[f(x)+4\right]^2}+\sqrt[3]{f(x)+4}+1\right]}\\
			&=&\lim\limits_{x\to 4}\left[\dfrac{f(x)+3}{x-4}\cdot  \dfrac{1}{\sqrt[3]{\left[f(x)+4\right]^2}+\sqrt[3]{f(x)+4}+1}\right]\\
			&=&1\cdot \dfrac{1}{3}=\dfrac{1}{3}. 
		\end{eqnarray*}
	 }
\end{ex}

% nguồn: Ok_TracNghiemLop11_Lan1
\begin{ex}%[1D3B2-1]%[ma:7ebe4da8]
	Cho dãy số $\heva{& u_1 = 4 \\ & u_{n+1} = \dfrac{u_n(n+4)}{n+3}, n \ge 1}$. Công thức tổng quát của dãy  số là
	\choice 
		{$u_n = 2n+2, \forall  n \in \mathbb{N}^*$}
		{$u_n  = -n+5, \forall n \in \mathbb{N}^*$}
		{\True  $u_n = n + 3, \forall n \in \mathbb{N}^*$}
		{$u_n = 3n+1, \forall n \in \mathbb{N}^*$} 
	 \loigiai{
		Ta có $u_{2} =  \dfrac{u_1(1+4)}{1+3} = 5$, $u_{3} = \dfrac{u_2(2+4)}{2+3} = 6$, $u_{4} = \dfrac{u_3(3+4)}{3+3} = 7$. \\
		Dự đoán: $u_n  = n + 3$ (*). \\
		Ta chứng minh (*) bằng quy nạp. Trước hết, $u_1 = 4 = 1 + 3$. Vậy (*) đúng khi $n = 1$. \\ 
		Giả sử (*) đúng khi $n = k, k \in \mathbb{N}^*$. Khi đó  $u_k = k + 3$.
		 Ta có
		$$u_{k+1} = \dfrac{u_k(k+4)}{k+3} = \dfrac{(k+3)(k+4)}{k+3} = k + 4 = (k + 1) + 3  $$
		Vậy (*) đúng  khi $n = k + 1$. Theo nguyên lý quy nạp, ta có điều phải chứng minh. \\
		Kết luận: $u_n = n + 3, \forall n \in  \mathbb{N}^*$.
	 }
\end{ex}

% nguồn: Ok_TracNghiemLop11_Lan1
\begin{ex}%[1D3B2-1]%[ma:97c6e599]
	Cho  dãy số $\left(u_n\right)$ biết $\heva{
	& u_1=2 \\ 
	& {u}_{n+1}=2u_n, \forall n  \in \mathbb{N}^{*}}$. Tìm số hạng tổng quát của dãy số này.
	 \choice
		{\True  $u_n=2^n$}
		 {$u_n={n}^{n-1}$}
		 {$u_n=2$}
		{$u_n={2}^{n+1}$} 
	 \loigiai{
	 Ta có
	$u_n=2{u}_{n-1}=2  \cdot 2{u}_{n-2}= \ldots ={2}^{n-1} \cdot u_1=2^n$.
	 }
\end{ex}

% nguồn: Ok_TracNghiemLop11_Lan1
\begin{ex}%[1D3B2-1]%[ma:34b22cc9]
	Dãy số có các số hạng cho bởi $0;\dfrac{1}{2};\dfrac{2}{3};\dfrac{3}{4};\dfrac{4}{5};\ldots $ có số hạng tổng quát  là công thức nào dưới đây?
	 \choice
		 {$u_n=\dfrac{n+1}{n}$}
		{$u_n=\dfrac{n}{n+1}$} 
		{\True $u_n=\dfrac{n-1}{n}$  }
		 {$u_n=\dfrac{n^2-n}{n+1}$}
	\loigiai{ 
	Đặt $ u_1=0 $, $ u_2=\dfrac{1}{2} $, $  u_3=\dfrac{2}{3} $, $ u_4=\dfrac{3}{4} $, $ u_5=\dfrac{4}{5}, \ldots $\\
	Ta có $ u_1=\dfrac{1-1}{1} $, $ u_2=\dfrac{2-1}{2} $, $ u_3=\dfrac{3-1}{3} $, $  u_4=\dfrac{4-1}{4} $, $ u_5=\dfrac{5-1}{5},\ldots $\\
	Vậy  $ u_n=\dfrac{n-1}{n}, \forall n\in \mathbb{N^*} $.
	} 
\end{ex}

% nguồn: Ok_TracNghiemLop11_Lan1
\begin{ex}%[1D3B2-1]%[ma:91eb9fdd]
	Cho dãy số  $(u_n)$ xác định bởi $ \heva{&u_1=2008 \\ &u_2=2009 \\ &2{{u}_{n+1}}=u_n+{{u}_{n+2}}\quad \text{ với } n\ge 1}$. Số hạng tổng quát của dãy là
	\choice 
		{\True $u_n=n+2007$} 
		{$u_n=2n+2006$} 
		 {$u_n=n-2007$}
		 {$u_n=-n+2008$}
	\loigiai{ 
		Ta có $  u_n-u_{n-1}=1\Rightarrow u_n=u_{n-1}+1$. Suy ra $$ u_n=(u_n-u_{n-1})+(u_{n-1}-u_{n-2})+\cdots +(u_2-u_1)+u_1=1+1+. . . +1+u_1=n-1+2008=n+2007.$$
	 }
\end{ex}

% nguồn: Ok_TracNghiemLop11_Lan1
\begin{ex}%[1D3B2-1]%[ma:6da5f84b]
	Cho dãy số $(u_n)$ với $u_1=6,u_n=u_{n-1}+5$. Khi đó $u_n$ được tính  theo biểu thức nào dưới đây?
	 \choice
		{\True  $u_n=5n+1$}
		 {$u_n=5(n+1)$}
		 {$u_n=5^n+1$}
		 {$u_n=5^{n+1}$}
	\loigiai{ 
		Vì  $u_1=6$ loại được các đáp án $u_n = 5(n + 1)$ và $u_n = 5^{n + 1}$.\\
		Theo hệ  thức truy hồi ta có $u_2=u_1 + 5= 6 + 5 = 11$ nên loại được đáp án $u_n = 5^n + 1$.\\
	Vậy đáp án $u_n = 5n + 1$ là phương  án cần tìm.}
\end{ex}

% nguồn: Ok_TracNghiemLop11_Lan1
\begin{ex}%[1D3K2-1]%[ma:6a8da0a5]
	Cho dãy $(u_n)$ được xác  định bởi $\heva{&u_1 = 1 \\ &u_{n+1} = u_n + n^2}$. Số hạng tổng quát của dãy số là
	\choice 
		{$u_n  = 1 + \dfrac{n(2n+1)(n+1)}{6}$}
		{$u_n  = 1 + \dfrac{(n-1)n(2n+2)}{6}$}
		{\True $u_n = 1  + \dfrac{(n-1)n(2n-1)}{6}$}
		{Tất cả  đều sai}
	\loigiai{ 
		Từ $u_{n+1} = u_n + n^2$, ta  có
		\begin{align*}
			&u_1 =  1\\
			&u_2 =  u_1 + 1^2\\
			&\vdots\qquad\vdots\\ 
			&u_n = u_{n-1} + (n-1)^2. 
		\end{align*}
		Cộng vế theo vế ta được $u_n = 1 + 1^2 + 2^2 + \cdots  + (n-1)^2$.\\
		Mặt khác ta có $1^2 +  2^2 + 3^2 + \cdots + (n-1)^2 + n^2 = \dfrac{n(n+1)(2n+1)}{6}$\\
		 $\Rightarrow 1^2+2^2+3^2+ \cdots + (n-1)^2=\dfrac{n(n-1)(2n-1)}{6}$\\
		 Vậy $u_n=1+\dfrac{(n-1)n(2n-1)}{6}.$
	 }
\end{ex}

% nguồn: Ok_TracNghiemLop11_Lan1
\begin{ex}%[1D3B2-1]%[ma:488a014a]
	Cho dãy số $(u_n)$ với $\heva{&u_1=1\\&u_{n+1}=u_n+(-1)^{2n}}$.  Số hạng tổng quát $u_n$ của dãy là
	 \choice
		{$u_n=2n-1$} 
		 {\True $u_n=n$}
		 {$u_n=n$}
		 {$u_n=n$}
	\loigiai{Ta có $u_{n+1}=u_n+(-1)^{2n}=u_n+1 \Rightarrow u_2=2;  u_3=3; u_3=4;...$ Dễ dàng đoán được $u_n=n$. \\
		Thật vậy, ta chứng minh được  $u_n=n\ (*)$ bằng phương pháp quy nạp như sau:
		 $(*)$ đúng khi $n=1$.\\
		Giả sử $(*)$ đúng với $n=k, k\in \mathbb{N}^*$ ta có $u_k=k$. Ta  đi chứng minh $(*)$ cũng đúng với $n=k+1$ tức là $u_{k+1}=k+1$. \\
	Thật vậy từ hệ thức  truy hồi ta có $u_{k+1}=u_k+(-1)^{2k}=k+1$ (đpcm). Vậy công thức số hạng tổng quát $u_n=n$.}
\end{ex}

% nguồn: Ok_TracNghiemLop11_Lan1
\begin{ex}%[1D3B2-1]%[ma:acb6cb3c]
	Cho dãy số $(u_n)$ xác định bởi $ \heva{&u_1=1  \\ &u_n=2u_{n-1}+3\quad\forall n\ge 2 }$. Số hạng tổng quát của dãy là
	 \choice
		{\True  $u_n=2^{n+1}-3 $}
		{$u_n=2^{n+2}-7 $} 
		{$u_n=2^{n}-1 $} 
		{$u_n=2^{n-1}+1  $}
	\loigiai{Số hạng  tổng quát của dãy là $u_n=2^{n+1}-3 $. Ta chứng minh bằng phương pháp quy nạp
		\begin{itemize}
			\item  Với $ n=1\Rightarrow u_1=2^{1+1}-3=1\Rightarrow $ bài toán đúng với $ n=1 $.
			\item Giả sử $u_k=2^{k+1}-3  $, ta chứng minh $u_{k+1}=2^{k+2}-3 $.\\
			Thật vậy, theo  công thức truy hồi ta có $u_{k+1}=2u_k+3=2(2^{k+1}-3)+3=2^{k+2}-3 $.
		\end{itemize}
	 }
\end{ex}

% nguồn: Ok_TracNghiemLop11_Lan1
\begin{ex}%[1D3B2-1]%[ma:2030f9a5]
	Số hạng tổng quát của dãy số $1; \dfrac{1}{2}; \dfrac{1}{3}; \dfrac{1}{4}; \dfrac{1}{5};  \ldots$ là
	\choice 
		 {\True $u_n=\dfrac{1}{n}$}
		 {$u_n=\dfrac{1}{n+1}$}
		 {$u_n=\dfrac{1}{2n}$}
		 {$u_n=\dfrac{1}{2n+1}$}
	\loigiai{ 
	Số hạng tổng quát của  dãy số $\dfrac{1}{1};\dfrac{1}{2};\dfrac{1}{3};\dfrac{1}{4};\dfrac{1}{5};\ldots$ là $u_n=\dfrac{1}{n}$.
	} 
\end{ex}

% nguồn: Ok_TracNghiemLop11_Lan1
\begin{ex}%[1D3B2-1]%[ma:c3d845b7]
	Cho dãy số có các số hạng đầu là $0{,}1$; $0{,}01$; $0{,}001$; $0{,}0001$; $\ldots$. Số hạng tổng quát của dãy số này có dạng 
	 \choice
		{\True  $u_n=\underbrace{0,00 \ldots 01}_{n\ \text{chữ số}\ 0}$}
		{$u_n=\underbrace{0,00 \ldots 01}_{n-1 \ \text{chữ  số} \ 0}$}
		 {$u_n=\dfrac{1}{10^{n-1}}$}
		 {$u_n=\dfrac{1}{10^{n+1}}$}
	 \loigiai{
		Ta  có\\
		Số  hạng thứ $1$ có $1$ chữ số $0$.\\
		Số hạng thứ $2$ có $2$ chữ  số $0$.\\
		Số hạng thứ $3$ có $3$  chữ số $0$.\\
		 …………………………….\\
	Suy  ra $u_n$ có $n$ chữ số $0$.}
\end{ex}

% nguồn: Ok_TracNghiemLop11_Lan1
\begin{ex}%[1D3K2-1]%[ma:66f1dc28]
	Cho dãy số $(u_n)$ với $\heva{&u_1=5 \\& u_{n+1}=u_n+n}$. Số hạng tổng quát của  dãy số này là
	\choice 
		 {$u_n=\dfrac{(n-1)n}{2}$}
		 {\True $u_n=5+\dfrac{(n-1)n}{2}$}
		{$u_n=5+\dfrac{(n+1)n}{2}$} 
		{$u_n=5+\dfrac{(n+1)(n+2)}{2}$} 
	\loigiai{Từ tính chất  của dãy ta có
		\begin{align*}
			&u_1=5 \\ 
			&u_2=u_1+1\\ 
			&u_3=u_2+2\\ 
			 &\cdots=\cdots\\
			&u_n=u_{n-1}+(n-1) 
		\end{align*}
		 Suy ra
		\begin{align*}
			&  u_1+u_2+\cdots +u_n=5+(u_1+u_2+\cdots +u_{n-1})+[1+2+\cdots +(n-1)]\\
			\Leftrightarrow \; & u_n=5+\dfrac{(n-1)n}{2}. 
		\end{align*}
	 }
\end{ex}

% nguồn: Ok_TracNghiemLop11_Lan1
\begin{ex}%[1D3B2-1]%[ma:6988b980]
	Cho $S_n=\dfrac{1}{1 \cdot  3}+\dfrac{1}{3 \cdot 5}+ \cdots +\dfrac{1}{(2n-1)(2n+1)}$ với $n \in \mathbb{N}^*$. Mệnh đề nào sau đây là đúng ?
	\choice 
		 {$S_n=\dfrac{n+1}{n}$}
		 {\True $S_n=\dfrac{n}{2n+1}$}
		{$S_n=\dfrac{n+2}{2n+7}$} 
		 {$S_n=\dfrac{n-1}{2n-1}$}
	\loigiai{ 
		Ta có $\dfrac{1}{n(n+2)}=\dfrac{1}{2} \cdot \dfrac{n+2-n}{n(n+2)}=\dfrac{1}{2}\left(\dfrac{n+2}{n(n+2)}-\dfrac{n}{n(n+2)}\right)=\dfrac{1}{2}\left(\dfrac{1}{n}-\dfrac{1}{n+2}\right).$ \\ 
		Do  đó
		\allowdisplaybreaks 
		\begin{eqnarray*}
			S_n&=&\dfrac{1}{1 \cdot 3}+\dfrac{1}{3 \cdot 5}+ \cdots +\dfrac{1}{(2n-1)(2n+1)}\\ 
			 &=&\dfrac{1}{2}\left(\dfrac{1}{1}-\dfrac{1}{3}+\dfrac{1}{3}-\dfrac{1}{5}+ \cdots +\dfrac{1}{2n-1}-\dfrac{1}{2n+1}\right)\\
			&=&\dfrac{1}{2}\left(\dfrac{1}{1}-\dfrac{1}{2n+1}\right)\\ 
			&=&\dfrac{1}{2} \cdot \dfrac{2n}{2n+1}\\ 
			 &=&\dfrac{n}{2n+1}.
		\end{eqnarray*}
		Vậy $S_n=\dfrac{n}{2n+1}$. 
	 }
\end{ex}

% nguồn: Ok_TracNghiemLop11_Lan1
\begin{ex}%[1D3B2-1]%[ma:edf4f288]
	Cho dãy số $(u_n)$ có $\heva{&u_1=3  \\&u_{n+1}=u_n+5 \text{ với }n\ge1}$. Số hạng tổng quát của dãy số là
	 \choice
		{$u_n=7n-4$} 
		{$u_n=4n-1$} 
		 {$u_n=n+2$}
		{\True $u_n=5n-2$} 
	 \loigiai{Ta có $u_2=8=5\cdot2-2, u_3=13=5\cdot3-2, u_4=18=5\cdot4-2 $.\\
	Dự đoán $\forall n\ge1$, ta có $u_n=5n-2$ $(1)$.\\ 
	Ta  chứng minh bằng phương pháp quy nạp.\\
	Với $n=1$, ta  có $u_1=3=5\cdot 1-2$. Vậy $(1)$ đúng với $n=1$.\\
	Giả sử $(1)$ đúng với  $n=k$ tức là ta có $u_k=5\cdot k-2$.\\
	Ta chứng minh $(1)$ cũng đúng với $n=k+1$,  tức là phải chứng minh
	$u_{k+1}=5.(k+1)-2$. \\ 
	Thật vậy, theo giả thiết quy nạp ta có $u_{k+1}=u_{k}+5=(5k-2)+5=5k+5-2=5(k+1)-2$.\\ 
	Vậy $(1)$ đúng với $n=k+1$  do đó $(1)$ đúng với mọi $n\ge1$.}
\end{ex}


%%%%%%%%%%%%%%%% ĐÚNG/SAI %%%%%%%%%%%%%%%%

% nguồn: Ok_dungsai_lop11_Lan2
\begin{ex}%[1D3H2-1]%[ma:f00e7bc1]
	Giả  sử ta có
	$\lim\limits_{x \to +\infty}f(x)=a$ và $\lim\limits_{x  \to +\infty}g(x)=b, a, b \in \mathbb{R}$.
	\choiceTF 
		{\True $\lim\limits_{x \to +\infty}[f(x) \cdot g(x)]=a  \cdot b$}
		{\True $\lim\limits_{x  \to +\infty}[f(x)-g(x)]=a-b$}
		{$\lim\limits_{x  \to +\infty}\dfrac{f(x)}{g(x)}=\dfrac{a}{b}$}
		{\True $\lim\limits_{x \to +\infty}[f(x)+g(x)]=a+b$} 
	\loigiai{ 
		\begin{itemchoice}
			\itemch \textbf{Đúng.} Vì $\lim\limits_{x  \to +\infty}[f(x) \cdot g(x)]=\lim\limits_{x \to +\infty}f(x) \cdot \lim\limits_{x \to +\infty}g(x)=a \cdot b$.
			\itemch \textbf{Đúng.} Vì $\lim\limits_{x \to +\infty}[f(x)-g(x)]=\lim\limits_{x \to +\infty}f(x)-\lim\limits_{x \to +\infty}g(x)=a-b$. 
			\itemch \textbf{Sai.} Vì  có thể $b=0$.
			 \itemch \textbf{Đúng.} Vì $\lim\limits_{x \to +\infty}[f(x)+g(x)]=\lim\limits_{x \to +\infty}f(x)+\lim\limits_{x \to +\infty}g(x)=a+b$.
		\end{itemchoice}
	} 
\end{ex}

% nguồn: Ok_dungsai_lop11_Lan2
\begin{ex}%[1D3H2-1]%[ma:a30cff02]
	Giả sử ta  có $\lim\limits_{x \to x_0} f(x)=a$ và $\lim\limits_{x \to x_0} g(x)=b$. Với $a$, $b \in \mathbb{R}$.
	\choiceTF 
		{\True $\lim\limits_{x  \to x_0}[f(x) \cdot g(x)]=a \cdot b$}
		{\True $\lim\limits_{x \to x_0}[2f(x)-g(x)]=2a-b$} 
		{$\lim\limits_{x  \to x_0} \dfrac{f(x)}{g(x)}=\dfrac{a}{b}$}
		{$\lim\limits_{x \to  x_0}\left[\dfrac{f(x)+2g(x)}{f(x)}\right]=a+\dfrac{2b}{a}$ với $(a \neq 0)$}
	\loigiai{ 
		\begin{itemchoice}
			\itemch Vì theo  định lí về giới hạn hữu hạn.
			 \itemch Vì $\lim\limits_{x \to x_0}[cf(x)]=c \lim\limits_{x \to x_0} f(x)$ $(c \in \mathbb{R})$ nên $\lim\limits_{x \to x_0}[2f(x)-g(x)]=2a-b$.
			\itemch  Vì $b$ có thể bằng $0$.
			\itemch Vì $\lim\limits_{x \to x_0}\left[\dfrac{f(x)+2g(x)}{f(x)}\right]=\lim\limits_{x \to x_0}\left[1+\dfrac{2g(x)}{f(x)}\right]=1+\dfrac{2b}{a}$; $(a \neq 0)$. 
		\end{itemchoice}
	} 
\end{ex}

% nguồn: Ok_dungsai_lop11_Lan2
\begin{ex}%[1D3H2-1]%[ma:fe1491ee]
	Cho  $\lim\limits_{x \to 3}f(x)=-2$.
	 \choiceTF
		{Giới hạn $\lim\limits_{x \to  3}\dfrac{x-3}{f(x)}$ bằng vô cực}
		{Giới hạn $\lim\limits_{x \to 3}\dfrac{f(x)}{x-3}$ tồn tại hữu hạn} 
		{\True Giới hạn $\lim\limits_{x  \to 3}\dfrac{f(x)}{(x-3)^2}=-\infty$}
		{\True Giới  hạn $\lim\limits_{x \to 3}\dfrac{(x-3)^2}{f(x)}=0$}
	\loigiai{ 
		\begin{itemchoice}
			\itemch  \textbf{Sai.} Ta có $\lim\limits_{x \to 3}\dfrac{x-3}{f(x)}=\dfrac{3-3}{-2}=0$.
			\itemch \textbf{Sai.} Ta có $\lim\limits_{x \to 3^{+}}\dfrac{f(x)}{x-3}=-\infty$ và  $\lim\limits_{x \to 3^{-}}\dfrac{f(x)}{x-3}=+\infty$ nên $\lim\limits_{x \to 3}\dfrac{f(x)}{x-3}$ không tồn tại.
			\itemch \textbf{Đúng.} Ta có $\lim\limits_{x \to 3}f(x)=-2<0$; $\lim\limits_{x  \to 3}(x-3)^2=0$ và $(x-3)^2>0$ khi $x \to 3$.
			 \\
			Vậy $\lim\limits_{x  \to 3}\dfrac{f(x)}{(x-3)^2}=-\infty$.
			\itemch \textbf{Đúng.} Ta có  $\lim\limits_{x \to 3}\dfrac{(x-3)^2}{f(x)}=\dfrac{(3-3)^2}{-2}=0$.
		\end{itemchoice}
	 }
\end{ex}


%%%%%%%%%%%%%%%% TỰ LUẬN %%%%%%%%%%%%%%%%

% nguồn: Ok_TuanLuan_Toan11_Lan1
\begin{bt}%[1D3Y2-1]%[ma:cf9c1e74]
	Cho dãy số $\left(u_n\right)$ viết dưới dạng khai triển $\dfrac{1}{2},\dfrac{2}{3},\dfrac{3}{4},\dfrac{4}{5},\dfrac{5}{6},\ldots$ Tìm  số hạng tổng quát $u_n$ của dãy số $\left(n\in \mathbb{N}^{\ast}\right)$.
	 \loigiai{
		Do $n\in \mathbb{N}^{\ast}$ nên  suy ra $u_n=\dfrac{n}{n+1}.$
	 }
\end{bt}

% nguồn: Ok_TuanLuan_Toan11_Lan1
\begin{bt}%[1D3Y2-1]%[ma:ccd6eaec]
	Cho dãy số $\left(u_n\right)$ viết dưới dạng khai triển $1,\dfrac{1}{4},\dfrac{1}{9},\dfrac{1}{16},\dfrac{1}{25},\ldots$. Tìm số hạng tổng  quát $u_n$ của dãy số $\left(n\in \mathbb{N}^{\ast}\right)$.
	 \loigiai{
		Do $n\in \mathbb{N}^{\ast}$ nên  suy ra $u_n=\dfrac{1}{n^2}.$
	} 
\end{bt}

% nguồn: Ok_TuanLuan_Toan11_Lan1
\begin{bt}%[1D3Y2-1]%[ma:cc8169d3]
	Cho dãy số $(u_n)$  với các số hạng đầu là $5;10;15;20;25;\ldots$. Tìm số hạng tổng quát của dãy số này.
	\loigiai{Từ tính chất các số hạng của dãy suy ra số hạng tổng quát của dãy $(u_n)$  là $u_n=5n$.}
\end{bt}

% nguồn: Ok_TuanLuan_Toan11_Lan1
\begin{bt}%[1D3Y2-1]%[ma:16cc3989]
	Cho dãy  số $\left(u_n\right)$ có số hạng tổng quát là $u_n=\dfrac{3n-2}{n+1}$, $\forall n\in\mathbb{N}$. Viết $\left( u_n\right)$ dưới dạng khai triển
	 \loigiai{
		Dạng tường minh  của $\left( u_n\right)$ là $-2;\dfrac{1}{2};\dfrac{7}{4};2;\dfrac{13}{6};\ldots$
	} 
\end{bt}

% nguồn: Ok_TuanLuan_Toan11_Lan2
\begin{ex}%[1D3Y2-1]%[ma:f4423230]
	Cho dãy số $\left(u_n\right)$ viết dưới dạng  khai triển $\dfrac{1}{2},\dfrac{2}{3},\dfrac{3}{4},\dfrac{4}{5},\dfrac{5}{6},\ldots$ Tìm số hạng tổng quát $u_n$ của dãy số $\left(n\in \mathbb{N}^{\ast}\right)$.
	 \loigiai{
		Do $n\in \mathbb{N}^{\ast}$ nên suy ra $u_n=\dfrac{n}{n+1}.$ 
	 }
\end{ex}

% nguồn: Ok_TuanLuan_Toan11_Lan2
\begin{ex}%[1D3Y2-1]%[ma:ce88db9f]
	 Cho dãy số $\left(u_n\right)$ viết dưới dạng khai triển $1,\dfrac{1}{4},\dfrac{1}{9},\dfrac{1}{16},\dfrac{1}{25},\ldots$. Tìm số hạng tổng quát $u_n$ của dãy số $\left(n\in \mathbb{N}^{\ast}\right)$.
	 \loigiai{
		Do $n\in \mathbb{N}^{\ast}$ nên  suy ra $u_n=\dfrac{1}{n^2}.$
	 }
\end{ex}

% nguồn: Ok_TuanLuan_Toan11_Lan2
\begin{ex}%[1D3Y2-1]%[ma:ce64d18e]
	Cho dãy số $(u_n)$ với các số hạng đầu là $5;10;15;20;25;\ldots$. Tìm số hạng tổng quát của dãy số này. 
	\loigiai{Từ tính chất các số hạng của dãy suy ra số hạng tổng quát của dãy $(u_n)$ là $u_n=5n$.} 
\end{ex}

% nguồn: Ok_TuanLuan_Toan11_Lan2
\begin{ex}%[1D3Y2-1]%[ma:6e5c20f7]
	Cho dãy số $\left(u_n\right)$ có số hạng tổng quát là $u_n=\dfrac{3n-2}{n+1}$, $\forall n\in\mathbb{N}$. Viết $\left( u_n\right)$  dưới dạng khai triển
	\loigiai{ 
		Dạng  tường minh của $\left( u_n\right)$ là $-2;\dfrac{1}{2};\dfrac{7}{4};2;\dfrac{13}{6};\ldots$
	 }
\end{ex}
